\uuid{IHqj}
\titre{Base et dimension}
\auteur{hanani ; mzari}
\organisation{exo7}
\licence{CC BY-NC-SA 4.0}
\datecreate{2020-01-20}
\niveau{L1}
\module{Algèbre}
\chapitre{Espace vectoriel}
\sousChapitre{Base}
\difficulte{3}
\theme{Base et dimension}
\exo7qcmid{lille-2.espaces-vectoriels.046}

\contenu{
\begin{question}
\qclassification{106.04}
\qtags{motcle=Bases de $\Cc$}
Quelles sont les assertions vraies ?
\begin{answers}
\bad{$\{1\}$ est une base de $\Cc$ comme $\Rr$-espace vectoriel}.
\good{$\{i\}$ est une base de $\Cc$ comme $\Cc$-espace vectoriel}.
\good{$\{i, 1+i\}$ est une base de $\Cc$ comme $\Rr$-espace vectoriel}.
\bad{$1$ et $i$ sont $\Cc$ linéairement indépendants}.
\end{answers}
\begin{explanations}  $\Cc$ est un $\Cc$-espace vectoriel de dimension $1$ et c'est un $\Rr$-espace vectoriel de dimension $2$. Par conséquent, pour tout $\alpha \in \Cc^*$, $\{\alpha\}$ est une base de $\Cc$ comme $\Cc$-espace vectoriel et  pour tous $\alpha, \beta \in \Cc^*$ tels que $\frac{\alpha}{\beta} \notin \Rr$,  $\{\alpha, \beta\}$ est une base de $\Cc$ comme $\Rr$-espace vectoriel.
\end{explanations}
\end{question}
}
